% Preamble for "Section 1: Sets and Binomial Expansions"
\usepackage[a4paper,margin=2.5cm]{geometry}
\usepackage{amsmath,amssymb,amsthm}
\usepackage{graphicx}
\usepackage{tikz}
\usetikzlibrary{calc,positioning,patterns}
\usepackage{enumitem}
\usepackage{xcolor}
\usepackage{tcolorbox}
\tcbuselibrary{skins,breakable}
\usepackage{fancyhdr}
\usepackage{hyperref}
\hypersetup{colorlinks=true,linkcolor=blue,urlcolor=blue,citecolor=blue}

\setlist[enumerate,1]{label=\arabic*.}
\setlist[enumerate,2]{label=(\alph*)}
\setlist[enumerate,3]{label=\roman*.}

% Colours matching a typical textbook palette
\definecolor{sectionblue}{HTML}{1F4E79}
\definecolor{boxblue}{HTML}{E8F1FA}
\definecolor{boxorange}{HTML}{FCEEE3}

\pagestyle{fancy}
\fancyhf{}
\fancyhead[L]{\small SECTION 1 \quad SETS AND BINOMIAL EXPANSIONS}
\fancyfoot[C]{\thepage}
\renewcommand{\headrulewidth}{0.4pt}

% Key Ideas box
\newtcolorbox{keyideas}{
  colback=boxblue, colframe=sectionblue, breakable,
  title={KEY IDEAS}, fonttitle=\bfseries
}

% Activity box
\newtcolorbox{activity}[1]{
  colback=boxorange, colframe=sectionblue, breakable,
  title={#1}, fonttitle=\bfseries
}

% Example box
\newtcolorbox{example}[1]{
  colback=white, colframe=sectionblue, breakable,
  title={#1}, fonttitle=\bfseries
}

\newcommand{\Union}{\cup}
\newcommand{\Inter}{\cap}

% ---------------------------------------------------------------------
% Reusable Venn-diagram backdrops.
% Each draws the outer universal-set box + three set outlines + labels;
% callers add \node's for the region contents appropriate to that stage,
% which lets a worked example build the same diagram up in several steps.
% ---------------------------------------------------------------------

% Three overlapping ellipses: two on top (black, blue), one on the
% bottom (red) overlapping both. Used for Examples 1.5, 1.7, 1.8.
\newcommand{\vennThreeEllipse}[4]{
  \draw (0,0) rectangle (9,6.3);
  \node at (8.7,6.0) {#4};
  \draw[black, thick] (3.0,4.0) ellipse (2.1 and 1.5);
  \node at (0.7,5.7) {#1};
  \draw[blue, thick] (5.5,4.0) ellipse (2.1 and 1.5);
  \node[blue] at (7.5,5.7) {#2};
  \draw[red, thick] (4.25,2.2) ellipse (2.4 and 1.5);
  \node[red] at (7.3,0.9) {#3};
}

% Two overlapping ellipses on top (black, blue), wide rectangle along
% the bottom (red) overlapping both well into their bodies (not just
% grazing the edge, so the triple-overlap region is large enough to
% label). Used for Example 1.6.
\newcommand{\vennEllipseRectBottom}[4]{
  \draw (0,0) rectangle (9,7.4);
  \node at (8.7,7.1) {#4};
  \draw[black, thick] (3.3,4.6) ellipse (2.1 and 2.2);
  \node at (1.0,7.05) {#1};
  \draw[blue, thick] (5.3,4.6) ellipse (2.1 and 2.2);
  \node[blue] at (7.6,7.05) {#2};
  \draw[red, thick] (1.0,0.8) rectangle (8.0,3.8);
  \node[red] at (0.5,2.2) {#3};
}

% Wide rectangle along the top (red), two overlapping ellipses below
% (blue, black) poking well up into it. Used for Activity 1.1.
\newcommand{\vennEllipseRectTop}[4]{
  \draw (0,0) rectangle (9,7.4);
  \node at (8.7,7.1) {#4};
  \draw[red, thick] (1.0,3.6) rectangle (8.0,6.6);
  \node[red] at (7.3,6.9) {#3};
  \draw[blue, thick] (3.3,2.8) ellipse (2.1 and 2.2);
  \node[blue] at (0.6,0.3) {#1};
  \draw[black, thick] (5.3,2.8) ellipse (2.1 and 2.2);
  \node[black] at (8.0,0.3) {#2};
}
